\documentclass[../../main.tex]{subfiles}
\begin{document}

\begin{center}
  {\fontsize{48pt}{58pt}\selectfont \textbf{2003}\normalsize\textbf{年度}}
\end{center}

\begin{table}[htb]
  \centering
  \caption{出題テーマと難易度}
  \begin{tabular}{|c|p{9cm}|c|}
    \hline
    \textbf{問題番号} & \textbf{テーマ} & \textbf{難易度} \\
    \hline
    第1問 & 3次関数のグラフと直線が囲む面積の比 & \\
    \hline
    第2問 & 長方形からの正方形の切り出し操作（連分数） & \\
    \hline
    第3問 & 三角形の内分点とベクトルによる領域の面積 & \\
    \hline
    第4問 & 関数列の漸化式と高次導関数の値 & \\
    \hline
  \end{tabular}
\end{table}

\begin{table}[htb]
  \centering
  \caption{解答}
  \renewcommand{\arraystretch}{1.4}
  \begin{tabular}{|c|c|p{9cm}|}
    \hline
    \textbf{問題番号} & \textbf{小問} & \textbf{答え} \\
    \hline
    第1問 & (1) & $y=bx$，$y=\left(\dfrac{1}{4}a^2+b\right)x$ \\
    \cline{2-3}
    & (2) & $S_1:S_2=1:16$ \\
    \hline
    第2問 & (1) & $a=\dfrac{4}{3},\dfrac{5}{3},\dfrac{5}{2},4$ \\
    \cline{2-3}
    & (2) & 最大値 $n+1$，最小値 $1+\dfrac{1}{n}$ \\
    \hline
    第3問 & - & $S=\dfrac{1}{3}-\dfrac{2}{9}\log2$ \\
    \hline
    第4問 & (1) & $x^{n+1}$の係数は$\dfrac{(n!)^2}{n}$ \\
    \cline{2-3}
    & (2) & $f_n^{(1)}(0)=0$，$f_n^{(2)}(0)=2$，$f_n^{(3)}(0)=12(n-1)$，$f_n^{(4)}(0)=144(n-1)(n-2)$ \\
    \hline
  \end{tabular}
\end{table}

\paragraph{試験形態}
4問．試験時間は要確認．

\end{document}
